\documentclass{article}
\title{Brackets and Grouping}
\author{ }
\date{}

\begin{document}
\maketitle

\section*{1. Basic brackets}
You can type ordinary round and square brackets directly.

\textbf{LaTeX source code:}
\begin{verbatim}
\[
(a+b) \qquad [x+y]
\]
\end{verbatim}
\textbf{Expected compiled output:}
\[
(a+b) \qquad [x+y]
\]

\section*{2. Automatically sized brackets}
Large expressions may need larger brackets.

\textbf{LaTeX source code:}
\begin{verbatim}
\[
\left(\frac{a}{b}\right)
\]
\end{verbatim}
\textbf{Expected compiled output:}
\[
\left(\frac{a}{b}\right)
\]
\textbf{Explanation:} \verb|\left| and \verb|\right| make brackets grow to fit the expression between them. They must be used as a matching pair.

\section*{3. Other bracket types}
\textbf{LaTeX source code:}
\begin{verbatim}
\[
\left[\frac{x+1}{x-1}\right]
\]
\end{verbatim}
\textbf{Expected compiled output:}
\[
\left[\frac{x+1}{x-1}\right]
\]

\section*{4. Practice}
Put large round brackets around a fraction and large square brackets around a sum.
\end{document}
